% Section 8.2, from lectures/L08/README.md: the goal, the evaluation question, and the next
% lecture. The README's note on simulating the watchdog is in the chapter's opener.

\needspace{12\baselineskip}
\section{Review}
\label{c8:sec:review}

\needspace{6\baselineskip}
\subsection*{What you should be able to do now}
After this chapter, you should be able to:
\begin{itemize}
  \item Explain what the watchdog timer protects against, and enable it, reset it, and observe
    what happens when it isn't reset in time.
\end{itemize}

\needspace{6\baselineskip}
\subsection*{Questions to test yourself}
\begin{enumerate}
  \item What is the watchdog timer for, and what's the practical difference between a program that
    resets it every loop iteration and one that never resets it at all?
\end{enumerate}

\needspace{6\baselineskip}
\subsection*{Looking ahead}
Now that EEPROM and the watchdog have shown what plain functions look like for a peripheral with
exactly one instance, Chapter~\ref{c9:ch} covers what changes once a program needs several
instances of the same kind of peripheral: structs as drivers.
